\section{Features and Leakage Control}
\label{sec:features}

The ten features are trailing ratios: returns over 1, 5 and 20 sessions; gaps to
5- and 20-session moving averages; realised volatility over 10 and 30 sessions;
intra-bar range; a volume anomaly; and day of week. Each uses bars up to and
including $t$, and rolling statistics require a full window.

Features are \emph{scale free} rather than price levels, which are non-stationary
and would let a model look accurate by copying today's price. The
reported quantity is error on the next close \emph{price}, while models are fitted
on the price \emph{change} and reconstructed as
$\hat{C}_{t+1} = C_t + \widehat{\Delta}_t$. The two are information-equivalent
because $C_t$ is known at decision time.

\paragraph{Boundary purge.}
With $y_t = C_{t+1}$, the last training label is a validation price and the last
validation label a sealed test price, so one row is dropped at each boundary. An
independent audit showed that, without it, perturbing one test price moved the
headline metric by 0.408 (Table~\ref{tab:splits}).

% GENERATED by scripts/build_latex_assets.py -- do not edit by hand.
\begin{table}[t]
\centering
\caption{Chronological split boundaries in the asset's own calendar. One row is purged at each boundary so that no label crosses a split: with $y_t = C_{t+1}$, the last training label would otherwise be a validation price and the last validation label a sealed test price. The test block remains unopened.}
\label{tab:splits}
\small
\setlength{\tabcolsep}{4pt}
\begin{tabular}{@{}l l l l r@{}}
\toprule
Asset & Training & Validation & Sealed test & Total \\
\midrule
AAPL & 2021-12-10..2024-11-08 & 2024-11-12..2025-11-03 & 2025-11-05..2026-10-27 & 1,224 \\
BTC & 2021-11-27..2024-11-07 & 2024-11-09..2025-11-01 & 2025-11-03..2026-10-28 & 1,797 \\
\bottomrule
\end{tabular}
\par\smallskip\footnotesize
Total is the number of usable feature rows after the trailing-window warm-up and the terminal unlabelled row are removed.
\end{table}


\paragraph{Enforcement.}
One test recomputes every feature from raw data with trailing formulas written
separately in the test, so a forward peek mismatches; another appends absurd future
prices and asserts earlier values are unchanged. Six canonical defects --- a
one-bar forward shift, an obfuscated equivalent, a centred window, a backward
fill, a whole-series statistic, and preprocessing fitted on train-plus-test ---
are then injected one at a time, and each must be detected.
